\documentclass[11pt]{article}
\usepackage[margin=1in]{geometry}
\usepackage{amsmath, amssymb}
\usepackage{bm}
\usepackage{graphicx}
\usepackage{booktabs}
\usepackage{hyperref}
\usepackage{enumitem}
\setlist[itemize]{itemsep=2pt, topsep=3pt}

\title{Applying Attribution-based Parameter Decomposition to the L$^4$-trained CiS toy model}
\author{Francisco Ferreira da Silva\\ \small (autonomous run, Claude Opus 4.7)}
\date{2026-04-17}

\newcommand{\Wein}{W_{\mathrm{in}}}
\newcommand{\Weout}{W_{\mathrm{out}}}

\begin{document}
\maketitle

\begin{abstract}
\noindent
We apply Attribution-based Parameter Decomposition (APD) \cite{bushnaqspd}
to our $L^4$-trained CiS toy models, both plain
($y = \Weout\,\mathrm{ReLU}(\Wein x)$) and with a random Gaussian embedding.
The APD paper shows APD recovering per-feature ``mechanisms'' in a similar
compressed-computation architecture where each feature's target is distinct
($y_i = x_i + \mathrm{ReLU}(c_i\,x_i)$ with random $c_i$). On our $L^4$
models, where features have identical targets ($y_i = \mathrm{ReLU}(x_i)$)
and the learned mechanism is combinatorial ReLU coding, APD does
\emph{not} recover per-feature components: the input and output halves of
each APD component decouple onto disjoint index sets, so no single
parameter component implements a feature's read-and-write mechanism. As a
sanity check we retrain a paper-style target (with random per-feature
coefficients) using the same architecture and find APD does produce
partial per-feature alignment. We attribute the contrast to our network's
mechanism being \emph{per-codeword} rather than per-feature: multiple
features route through the same neurons, so there is no per-feature
circuit to recover.
\end{abstract}

\section{Setup}

\subsection{APD in brief}

APD \cite{bushnaqspd} trains a second, ``decomposed'' copy of a target
model whose weight matrices $W$ are replaced by a sum of $C$ rank-$m$
components:
\[
W = \sum_{c=1}^{C} A_c B_c,\qquad A_c \in \mathbb{R}^{d_\mathrm{in}\times m},\quad B_c \in \mathbb{R}^{m\times d_\mathrm{out}}.
\]
Training uses four losses:
\begin{itemize}
\item \textbf{param-match:} $\|\sum_c A_cB_c - W_\mathrm{target}\|^2$ per layer,
  so the full sum reconstructs the target's weight.
\item \textbf{topk-recon:} at each step, a gradient-attribution score is computed
  per component and a top-$k$ subset is kept; the decomposed model is run
  through only those components and its output is pushed to match the
  target's output.
\item \textbf{act-recon:} same as above but on the post-MLP-in activations
  (intermediate representations).
\item \textbf{Schatten regulariser:} encourages components to be low rank.
\end{itemize}
Together these push the decomposition to be sparse (few components active
per input), reconstructive (the sum equals the original weight), and
functional (the top-$k$ components can reproduce the target output).

\subsection{Target models}

We decompose five target models (all using our CiS architecture
$y = \Weout\,\mathrm{ReLU}(\Wein x)$ without residual connection and
trained under $L^4$ loss):
\begin{itemize}
\item \texttt{plain\_20f\_5n}: $F=20$ features, $n=5$ neurons.
\item \texttt{plain\_20f\_2n}: $F=20$, $n=2$.
\item \texttt{plain\_100f\_10n}: $F=100$, $n=10$.
\item \texttt{embed\_20f\_5n\_D80}: embedded variant with $D=80$, otherwise
  like \texttt{plain\_20f\_5n}.
\item \texttt{embed\_20f\_2n\_D40}: embedded variant with $D=40$, otherwise
  like \texttt{plain\_20f\_2n}.
\end{itemize}
As a \emph{pipeline sanity check}, we also train a ``paper-style'' target
that mimics the APD paper's toy-model task: target $y_i = \mathrm{ReLU}(c_i\,x_i)$
with random $c_i \in [1, 2]$, same architecture and size as
\texttt{plain\_20f\_5n}.

\subsection{Adapting the APD code to our architecture}

The APD paper uses a target architecture with a residual connection,
$y = W_U(\mathrm{resid} + \mathrm{MLP}(\mathrm{resid}))$. Our target has no
residual. We apply a minimal monkey-patch to \texttt{ResidualMLPModel} and
\texttt{ResidualMLPSPDModel} that replaces the residual add
(\texttt{resid $\leftarrow$ resid + MLP(resid)}) with a pure substitution
(\texttt{resid $\leftarrow$ MLP(resid)}). With this change we load our
trained weights as the target and set $W_E$, $W_U$ to identity for
\texttt{plain\_*} and to the fixed random $E$, $E^T$ for \texttt{embed\_*}.
We verified numerically that the patched APD model reproduces our
model's output to float precision (max difference $< 10^{-7}$ on 512
random sparse samples).

\subsection{APD training configuration}

Following the paper, we used:
\begin{itemize}
\item $C=30$ components for $F=20$ targets (20 features $<$ 30 components,
  small overhead for ``dying'' components), $C=130$ for $F=100$.
\item 10{,}000 steps, batch size 256, Adam lr $10^{-3}$, cosine schedule.
\item Gradient-based attribution, \texttt{batch\_topk=True}, \texttt{topk=1.28}
  (paper value; slightly above expected active features per sample).
\item \texttt{unit\_norm\_matrices=True}. Coefficients: param-match 1,
  topk-recon 1, act-recon 1, Schatten 10 with $p=0.9$.
\item Input distribution during APD training: $p = 0.05$ so that
  $E[\text{active per sample}] = 1$ (matching the paper's ratio of
  $\mathrm{topk}/E[\text{active}] \approx 1.28$). For $F=100$ we use
  $p=0.02$, $\mathrm{topk}=2.56$.
\end{itemize}

\section{Results}

\subsection{Pipeline sanity: paper-style target}

On the paper-style target (per-feature random coefficients), APD recovers
partial per-feature decomposition (Figure~\ref{fig:paperstyle}):
\begin{itemize}
\item 19/30 components on the input side are ``dead'' (essentially zero
  weight). Of the 11 alive, 7 are monosemantic (read strongly from a single
  feature).
\item On the output side, 9/30 dead, 9 monosemantic, 10 duo.
\item \textbf{2/30 components are alive on both sides with the same
  dominant feature} (the ``joint monosemantic'' criterion). These are
  clean per-feature mechanisms.
\item Scrubbing experiment: ablating half the components while keeping
  those whose dominant feature is active in the current sample gives
  mean MSE $7.7\times10^{-4}$; doing the opposite (anti-scrubbing) gives
  $6.2\times10^{-3}$. Ratio $\approx 8\times$. The decomposition is
  functionally meaningful.
\end{itemize}

\begin{figure}[h]
\centering
\includegraphics[width=\textwidth]{apd_decomposition/out/paper_style_20f_5n/attribution_heatmap.png}
\caption{\textbf{APD on paper-style target.} Left: per-feature read strength
of each input-side component (rows are components sorted by dominant
feature). Middle: per-feature write strength of each output-side component.
Right: component category counts, in-side vs out-side. Some components
have a clear dominant feature on both sides, matching the paper's
qualitative finding.}
\label{fig:paperstyle}
\end{figure}

This is worse than the paper's reported result ($\sim 100/100$ features
covered) but consistent with our smaller, more compressed task
($F/n = 4$ here vs $2$ in the paper). It confirms that our APD pipeline
is working and can in principle find per-feature components.

\subsection{Main finding: APD on $L^4$ CiS models gives polysemantic
components with decoupled in/out halves}

On our $L^4$-trained plain and embedded models the picture is very
different. Figures~\ref{fig:plain20f5n}--\ref{fig:plain20f2n} show
heatmaps from \texttt{plain\_20f\_5n} and \texttt{plain\_20f\_2n}:
\begin{itemize}
\item Many components are alive on the input side and many on the output
  side, but the alive sets are \emph{disjoint}: almost no component has
  both halves alive simultaneously. For \texttt{plain\_20f\_5n} we get
  15 alive on IN, 15 alive on OUT, and $0/30$ alive on both.
\item Each component that is alive on its one side is polysemantic --
  reads or writes multiple features. No component cleanly maps to a
  single feature's mechanism.
\item The scrubbing/anti-scrubbing ratio is $< 1$: the ``dominant
  feature'' of a component is essentially random, so keeping
  ``relevant'' components doesn't preserve behaviour.
\end{itemize}

Table~\ref{tab:summary} summarises across all runs.

\begin{table}[h]
\centering
\caption{APD decomposition summary. ``alive both'' = components alive on
both input and output sides. ``aligned'' = of those, how many have the
same dominant feature on both sides. ``anti/scrub'' is the ratio of
anti-scrubbed MSE to scrubbed MSE; $>1$ means the decomposition is
behaviourally meaningful.}
\label{tab:summary}
\small
\begin{tabular}{lrrrrrrr}
\toprule
target & F & C & $\mathrm{alive\_both}$ & aligned & joint mono & anti/scrub & MSE full \\
\midrule
paper\_style\_20f\_5n    & 20  & 30  & 2  & 2  & 2 & \textbf{8.05} & $3.5\times10^{-4}$ \\
\midrule
plain\_20f\_5n           & 20  & 30  & 0  & 0  & 1 & 0.39 & $2.3\times10^{-4}$ \\
plain\_20f\_2n           & 20  & 30  & 0  & 0  & 1 & 0.19 & $2.1\times10^{-4}$ \\
plain\_100f\_10n         & 100 & 130 & 9  & 5  & 2 & 0.39 & $1.0\times10^{-5}$ \\
embed\_20f\_5n\_D80      & 20  & 30  & 5  & 2  & 2 & 0.75 & $7.0\times10^{-5}$ \\
embed\_20f\_2n\_D40      & 20  & 30  & 0  & 0  & 3 & 0.21 & $7.0\times10^{-5}$ \\
\bottomrule
\end{tabular}
\end{table}

\begin{figure}[h]
\centering
\includegraphics[width=\textwidth]{apd_decomposition/out/plain_20f_5n/attribution_heatmap.png}
\caption{\textbf{APD on \texttt{plain\_20f\_5n}.} Both IN and OUT sides
have many alive components, each covering multiple features (dense
horizontal bands). The alive components on each side are at different
indices: no component has a coherent ``read feature $j$ then write feature
$j$'' structure. Right panel: no monosemantic or duo components.}
\label{fig:plain20f5n}
\end{figure}

\begin{figure}[h]
\centering
\includegraphics[width=\textwidth]{apd_decomposition/out/plain_20f_2n/attribution_heatmap.png}
\caption{\textbf{APD on \texttt{plain\_20f\_2n}.} More extreme version of
the same pattern. With only 3 codewords for 20 features, the OUT-side
alive components each write to (nearly) all 20 features simultaneously --
a direct signature of codeword sharing.}
\label{fig:plain20f2n}
\end{figure}

\subsection{Why? Codeword sharing means no per-feature circuit}

Our $L^4$-trained CiS model implements \emph{combinatorial ReLU coding}
(see companion report): features are routed to neurons via sign patterns
of $\Wein$ columns (``codewords''), and many features can share a single
codeword. When features $i$ and $j$ share a codeword, they pass through
the exact same neurons, so the ``mechanism'' of feature $i$ is identical
to the mechanism of feature $j$. There is no per-feature circuit to
recover.

APD's ansatz assumes features have dissociable mechanisms: each component
should implement one feature's input-to-output map. In the paper-style
target, random per-feature coefficients $c_i$ create genuinely distinct
per-feature mechanisms, so APD recovers them (partially). In our
\emph{uniform-target} setting ($y_i = \mathrm{ReLU}(x_i)$ for all $i$),
features share circuits, and APD has no way to separate them -- so it
returns a polysemantic decomposition whose input and output halves are
additionally decoupled.

The 20f/2n case makes this especially vivid: with only 3 codewords to
cover 20 features, the OUT-side component heatmap shows each alive
component writing to essentially all 20 features (wide horizontal
bands). This is the signature of shared circuitry on the write side.

\subsection{Plain vs embedded}

The embedded variants (\texttt{embed\_*}) show the same qualitative
picture as the plain ones (Table~\ref{tab:summary}). One small
quantitative difference: \texttt{embed\_20f\_5n\_D80} does have 5
components alive on both sides (vs 0 for the plain variant at the same
F,n), and 2 of those are in/out aligned. Scrubbing ratio improves from
0.39 to 0.75 (still $< 1$, still not a meaningful decomposition, but
visibly less pathological). Our reading: the random embedding breaks
some residual symmetry between input and output sides, which makes
pairing slightly easier, but doesn't change the underlying story.

\subsection{Scaling: \texttt{plain\_100f\_10n}}

Our largest target has 100 features, 10 neurons. APD with $C=130$ produces
121 joint-dead components and $\sim 9$ alive on both sides (5 aligned).
Two features get a clean monosemantic component. The overall picture
matches the $F=20$ runs: polysemantic, decoupled in/out halves, no
per-feature decomposition. Reconstruction MSE is $10^{-5}$ — better than
at $F=20$ because the model has more neurons per feature, so the
decomposition sum is easier.

\subsection{Anti/scrub ratio: a summary diagnostic}

Table~\ref{tab:summary} compresses a lot of information into one column:
the scrubbing ratio ``anti/scrub''. When this is $> 1$, APD has found a
meaningful association between components and features: ablating the
components attributed to active features (anti-scrub) does more damage
than ablating unrelated ones (scrub). The paper-style target sits at
8.05. Our five $L^4$ targets all sit between 0.19 and 0.75, i.e.~the
``dominant feature'' criterion over components is no better than random.
This is the single cleanest summary of the null result.

\section{Sanity-checks and caveats}

\begin{itemize}
\item \textbf{Reconstruction works.} For every run the full-decomposition
  MSE (sum of all $C$ components) is $\sim 10^{-4}$, essentially matching
  the target model's output. The decomposition isn't ``broken'' at the
  level of reconstructing weights; it just doesn't factorise into
  per-feature pieces.
\item \textbf{Hyperparameter sensitivity.} We tried three \texttt{topk}
  regimes (\texttt{topk=1.0} at $p=0.02$, $\mathrm{topk}=0.5$ at $p=0.02$,
  and the final \texttt{topk=1.28} at $p=0.05$) plus
  \texttt{data\_generation\_type="exactly\_one\_active"} and
  \texttt{exact\_topk=True}. Qualitatively all variants gave
  polysemantic, decoupled-in-out decompositions for our $L^4$ targets.
  The paper-style target gave partially meaningful decomposition in all
  tried configurations.
\item \textbf{Paper replicability.} Our 2/30 joint-monosemantic count on
  the paper-style target is lower than the paper's reported result. The
  difference likely comes from (a) tighter compression ($F/n=4$ vs $2$ in
  the paper), (b) no residual connection in our target (theirs can route
  some information around the MLP), and (c) a shorter, less-tuned run.
  We did not attempt to tune APD hyperparameters beyond matching the
  paper's config to our setup.
\item \textbf{Scrubbing caveat.} Our scrubbing criterion uses the
  ``dominant feature'' of each component (argmax over per-feature
  attribution). When components are polysemantic this criterion is noisy;
  the $<1$ ratios should be read as ``no meaningful dominant feature''
  rather than ``APD produces harmful ablations.''
\end{itemize}

\section{Conclusion}

APD applied to our $L^4$-trained CiS toy models does \emph{not} produce a
clean per-feature decomposition. Components' input and output halves end
up on disjoint sets of indices, and each alive component is polysemantic.
A pipeline sanity check on a paper-style target confirms APD can produce
at least partially per-feature components when the target has
distinguishable per-feature mechanisms.

We interpret this as consistent with the mechanism we identify in the
companion work: the $L^4$ model implements combinatorial ReLU coding,
where multiple features share the same neurons via shared codewords.
There is no per-feature circuit, so there is nothing per-feature for APD
to recover. The $L^2$-trained (``naive'') variants of the same
architecture would presumably decompose cleanly per-feature, since they
implement one feature per neuron; we didn't verify this here but it is
a natural next experiment.

\begin{thebibliography}{9}
\bibitem{bushnaqspd} Bushnaq \textit{et al.} (2025), \emph{Interpretability
in Parameter Space: Minimizing Mechanistic Description Length with
Attribution-based Parameter Decomposition}, arXiv:2501.14926.
\end{thebibliography}

\end{document}
